%-----------------------------------------------------
\section{Introduction}
%-----------------------------------------------------

Structures in service accumulate damage, and the traditional answer to that is periodic manual inspection.
Scheduled inspection has two weaknesses. It only sees the structure at fixed intervals, so damage that
initiates and grows between two inspections goes unnoticed, and it depends heavily on the judgement of the
inspector. Structural health monitoring replaces that with permanently installed sensors and an automated
diagnostic algorithm, so the condition of the structure can be interrogated whenever it is needed rather than
when the maintenance calendar allows.

The aerospace case is the one that motivated this work. Airframes are thin-walled, fatigue-critical, and
expensive to disassemble for inspection. A method that can interrogate a large area of skin from a small
number of permanently bonded transducers avoids most of that disassembly, and ultrasonic guided waves
are the main candidate for the role. In a plate whose thickness is comparable to the acoustic wavelength, the wave
is guided by the two free surfaces and travels tens of centimetres with modest attenuation, which means a
sparse network of transducers can cover a wide area. Any crack, delamination or corrosion patch along the
path between two transducers scatters the wave and changes the received signal.

\subsection{Scope and structure of this report}

The work reported here was carried out in three stages, and the report follows the same order.

\begin{enumerate}[leftmargin=*,itemsep=2pt]
  \item Section~\ref{sec:review} surveys the families of SHM methods and compares them against the
  requirements of a thin metallic aerospace panel, which leads to the choice of piezoelectric guided-wave
  monitoring. Section~\ref{sec:physics} then sets out the physical basis: the piezoelectric constitutive
  equations and Lamb-wave dispersion in a plate.

  \item Section~\ref{sec:stage1} builds a pitch-catch finite-element model in Abaqus with one actuator and
  one sensor, runs it in healthy and cracked states, and compares the two received signals using
  time-domain damage indices. The question it answers is whether damage is present.

  \item Section~\ref{sec:stage2} rebuilds the model with four transducers, giving six independent sensing
  paths, and combines the per-path damage indices into a two-dimensional damage-probability map. The
  question it answers is where the damage is.
\end{enumerate}

Sections~\ref{sec:stage1} and~\ref{sec:stage2} are reported as a single progression rather than as separate
exercises. Detection is established first, because a method that cannot detect damage has nothing to
localise, and the four-transducer model then extends the same finite-element approach with the
transducers, excitation design and crack representation that localisation requires. This report discusses
each of those additions where it arises.

%-----------------------------------------------------
\section{Structural health monitoring: levels, methods and selection}
\label{sec:review}
%-----------------------------------------------------

\subsection{Levels of diagnosis}

SHM tasks are usually organised into a hierarchy of increasing diagnostic ambition. The classification
introduced by Rytter defines four levels, and later work has added two more \cite{Rytter1993,Gharehbaghi2021}.

\begin{table}[H]
\centering
\caption{Diagnostic levels in structural health monitoring.}
\fitwidth{%
\begin{tabular}{c l p{7.4cm}}
\toprule
Level & Question answered & Typical output \\
\midrule
1 & Is there damage? & Binary or scalar flag; novelty detection against a threshold \\
2 & Where is it? & Coordinates, region or a damage-probability map \\
3 & How severe is it? & Crack length, stiffness loss, section-area reduction \\
4 & What is the remaining life? & Growth prediction, time to an unacceptable state \\
5 & What type of damage is it? & Crack, corrosion, delamination, loose joint \\
6 & What should be done? & Maintenance decision, load restriction \\
\bottomrule
\end{tabular}}
\label{tab:levels}
\end{table}

The distinction matters for this report because the three stages of work map onto the first two levels. The
two-transducer model of Section~\ref{sec:stage1} reaches Level~1, and the four-transducer network of
Section~\ref{sec:stage2} reaches Level~2. Levels 3 and above were not attempted.

\subsection{Survey of monitoring methods}

\paragraph{Vibration-based methods.}
The structure is excited at low frequency and its modal properties, natural frequencies, mode shapes, modal
curvature or flexibility, are tracked over time. Damage reduces local stiffness and shifts those properties
\cite{Yang2021Vib,Rabi2024}. The instrumentation is cheap and the method covers the whole structure at
once, which is why it dominates civil-infrastructure monitoring. The weakness is sensitivity: global modal
parameters respond weakly to small local damage, and they respond strongly to temperature, which makes
early crack detection in a metallic panel difficult.

\paragraph{Acoustic emission.}
Rather than sending a wave in, acoustic emission listens for the elastic energy released by the damage itself
as it grows \cite{Ghadarah2023,Muir2021}. It is passive and can catch damage at the moment of initiation,
which no active method can do. In exchange it only sees damage while it is actively growing, it needs
continuous listening at high sample rates, and it is vulnerable to operational noise.

\paragraph{Ultrasonic guided waves.}
A transducer launches a Lamb wave into the plate, and other transducers record what arrives
\cite{YangZ2023,Alem2016}. Because the wave travels far in a thin-walled structure, a small number of
transducers covers a large area, and because the excitation frequency is chosen by the operator the
wavelength can be matched to the size of the damage that must be found. The costs are a dispersive,
multi-modal wavefield that has to be understood before the signals can be interpreted, and a strong
dependence on a healthy-state baseline that itself drifts with temperature.

\paragraph{Fibre-optic sensing.}
Fibre Bragg gratings measure strain at discrete points along a single fibre, and Brillouin or Rayleigh
scattering methods give strain distributed continuously along its length
\cite{Alhussein2025,Barrias2016,Garcia2015}. The fibres are light, immune to electromagnetic interference
and can be embedded in a composite laminate. They measure strain, though, and strain is a local quantity:
damage away from the fibre produces little signal.

\paragraph{Electromechanical impedance.}
A bonded PZT patch is driven over a frequency sweep, typically above 30\,kHz, and its electrical impedance
is measured. Because the patch is coupled to the structure, the electrical impedance carries a copy of the
local mechanical impedance, and damage near the patch changes it \cite{Na2018}. The method is very
sensitive but only within a small radius of each patch, which makes area coverage expensive.

\paragraph{Vision-based and thermographic methods.}
Camera-based inspection detects surface cracks and measures displacement fields, and infrared thermography
detects subsurface defects through the thermal contrast they produce under active heating
\cite{He2020,Mardanshahi2025}. Both need line of sight and controlled conditions, so they suit periodic
inspection better than permanent monitoring.

\paragraph{Machine-learning approaches.}
Machine learning is a processing layer applied on top of any of the sensing methods above, to classify
damage, detect novelty, or fuse multiple data sources \cite{Jia2023,Hassani2023}. It becomes attractive once
a large body of labelled signals exists, which in practice usually means simulated signals.

\subsection{Selection}

Table~\ref{tab:compare} summarises the comparison against the requirements of a thin metallic panel:
sensitivity to a small crack, area covered per sensor, and suitability for permanent installation.

\begin{table}[H]
\centering
\caption{Comparison of the main SHM method families for thin metallic panels.}
\fitwidth{%
\begin{tabular}{l c c c}
\toprule
Method & Small-crack sensitivity & Area per sensor & Permanent install \\
\midrule
Vibration-based        & Low        & Whole structure & Yes \\
Acoustic emission      & High\,$^{*}$ & Large         & Yes \\
Guided waves (PZT)     & High       & Large           & Yes \\
Fibre optic (FBG)      & Medium     & Local           & Yes \\
Electromechanical impedance & High  & Very local      & Yes \\
Thermography / vision  & Medium     & Surface only    & No \\
\bottomrule
\end{tabular}}\\[4pt]
{\small $^{*}$ Only while the damage is actively growing.}
\label{tab:compare}
\end{table}

Guided-wave monitoring with bonded piezoelectric transducers is the only entry that combines high
sensitivity to a small crack with a large area per sensor and permanent installation. The same transducer can
act as actuator and as sensor, and the hardware is thin, light and cheap enough to be distributed over a
structure \cite{Qing2019,Ju2023}. That is the method used for the rest of this report.

%-----------------------------------------------------
\section{Physical basis}
\label{sec:physics}
%-----------------------------------------------------

\subsection{Piezoelectric constitutive equations}

A piezoelectric material couples its mechanical and electrical states. In the strain-charge form, the strain
tensor $\mathbf{S}$ and the electric displacement $\mathbf{D}$ are

\begin{align}
  \mathbf{S} &= \mathbf{s}^{E}\,\mathbf{T} + \mathbf{d}^{\mathsf{T}}\,\mathbf{E} \\
  \mathbf{D} &= \mathbf{d}\,\mathbf{T} + \boldsymbol{\varepsilon}^{T}\,\mathbf{E}
\end{align}

where $\mathbf{T}$ is stress, $\mathbf{E}$ the electric field, $\mathbf{s}^{E}$ the compliance at constant
field, $\boldsymbol{\varepsilon}^{T}$ the permittivity at constant stress, and $\mathbf{d}$ the piezoelectric
strain-charge coupling matrix. The two equations are what make a single patch usable in both directions.
Applying a voltage produces a strain, so the patch acts as an actuator; deforming the patch produces a charge,
so the patch acts as a sensor.

For a disc poled through its thickness, the components that matter are $d_{33}$, which couples the poling
direction to itself, $d_{31}=d_{32}$, which couples the in-plane directions to the poling direction and is
what converts an in-plane Lamb wave into a measurable voltage, and the shear terms $d_{15}=d_{24}$. The
values used throughout this work are those of APC850, listed in Table~\ref{tab:mat2}.

\subsection{Lamb waves in a thin plate}

In a plate bounded by two traction-free surfaces, the bulk longitudinal and shear waves reflect and interfere
until only certain guided modes survive. These are the Lamb modes, and they separate into symmetric and
antisymmetric families governed by the Rayleigh--Lamb equations,

\begin{align}
  \text{symmetric:}\qquad
  \frac{\tan(qh)}{\tan(ph)} &= -\frac{4k^{2}pq}{\left(q^{2}-k^{2}\right)^{2}} \\[4pt]
  \text{antisymmetric:}\qquad
  \frac{\tan(qh)}{\tan(ph)} &= -\frac{\left(q^{2}-k^{2}\right)^{2}}{4k^{2}pq}
\end{align}

with

\begin{equation}
  p^{2} = \left(\frac{\omega}{c_{L}}\right)^{2} - k^{2}, \qquad
  q^{2} = \left(\frac{\omega}{c_{T}}\right)^{2} - k^{2},
\end{equation}

where $h$ is the half-thickness, $k$ the wavenumber, $\omega$ the angular frequency, and $c_{L}$, $c_{T}$
the bulk longitudinal and shear velocities. These are transcendental in $k$ and have to be solved numerically
at each frequency. Doing so gives the phase velocity $c_{p}=\omega/k$ and the group velocity
$c_{g}=\mathrm{d}\omega/\mathrm{d}k$ for each mode.

Two consequences shape the modelling decisions that follow. First, the velocities depend on frequency, so a
pulse of finite bandwidth spreads as it travels and the received packet is wider than the transmitted one.
Second, at least two modes exist at any frequency, the fundamental symmetric $S_{0}$ and antisymmetric
$A_{0}$, and they travel at different speeds, so the sensor sees them as separate arrivals. Keeping the
excitation below the cut-on frequency of the higher modes keeps the picture to those two.

For the 1.6\,mm aluminium plate used in this work, driven at 200\,kHz, solving the equations above gives the
values in Table~\ref{tab:dispersion}.

\begin{table}[H]
\centering
\caption{Lamb-wave properties of the 1.6\,mm aluminium plate at 200\,kHz.}
\begin{tabular}{l c c c}
\toprule
Mode & Phase velocity (m/s) & Group velocity (m/s) & Wavelength (mm) \\
\midrule
$S_{0}$ & 5348 & 5332 & 26.7 \\
$A_{0}$ & 1580 & 2627 & 7.90 \\
\bottomrule
\end{tabular}
\label{tab:dispersion}
\end{table}

The $S_{0}$ mode is the faster of the two and arrives first, which is why every damage index in this report is
computed on the first arriving packet. The $A_{0}$ mode has the shorter wavelength, 7.90\,mm against
26.7\,mm, so it is the mode that sets the element size the mesh must resolve.

%-----------------------------------------------------
\section{Stage 1: crack detection with a two-transducer path}
\label{sec:stage1}
%-----------------------------------------------------

If a wave is launched from one bonded transducer and received by another, does a crack on the path change
the received signal enough to be measured? That is the only question the first model was built to answer.

\subsection{Model description}

An aluminium plate carries two PZT discs, one acting as actuator and one as sensor, placed on opposite faces
at either end of the plate so that the direct path between them passes over the centre. The crack sits at that
centre, directly across the path. Geometry is given in Table~\ref{tab:geom2} and material data in
Table~\ref{tab:mat2}.

\begin{table}[H]
\centering
\caption{Geometry of the two-transducer model.}
\begin{tabular}{l l}
\toprule
\multicolumn{2}{l}{\textbf{Aluminium plate}}\\
Dimensions & $177.8 \times 76.2 \times 1.6$\,mm \\
Centre & $(0,0)$\,mm \\
\midrule
\multicolumn{2}{l}{\textbf{Crack}}\\
Count & 1 \\
Dimensions & $1 \times 25 \times 1$\,mm \\
Centre & $(0,0)$\,mm \\
Notch root radius & 1\,mm \\
\midrule
\multicolumn{2}{l}{\textbf{PZT discs}}\\
Count & 2 \\
Dimensions & $r = 6.35$\,mm, $t = 0.254$\,mm \\
Actuator position & $(-69.85,\ 0,\ -1.6)$\,mm \\
Sensor position & $(+69.85,\ 0,\ +1.6)$\,mm \\
\bottomrule
\end{tabular}
\label{tab:geom2}
\end{table}

\begin{table}[H]
\centering
\caption{Material data for the plate and the piezoelectric discs.}
\begin{tabular}{l l l}
\toprule
Property & Aluminium plate & PZT disc \\
\midrule
Material          & Al 2024-T3 & APC850 \\
Density           & $2700$\,kg/m$^3$ & $7600$\,kg/m$^3$ \\
Young's modulus   & $69$\,GPa & $63$\,GPa \\
Poisson's ratio   & $0.33$ & $0.35$ \\
Permittivity      & --- & $1.6851\times10^{-8}$\,F/m \\
Charge constants  & --- & $d_{15}=d_{24}=590$\,pC/N \\
                  & --- & $d_{31}=d_{32}=-175$\,pC/N \\
                  & --- & $d_{33}=400$\,pC/N \\
\bottomrule
\end{tabular}\\[4pt]
{\small All remaining components of the charge-constant matrix are zero.}
\label{tab:mat2}
\end{table}

The distance between the two discs is $139.7$\,mm. With the $S_{0}$ group velocity of
Table~\ref{tab:dispersion}, the direct arrival is expected at $139.7/5332 = 26.2$\,\textmu s, which is worth
remembering when the signals are examined in Section~\ref{sec:stage1results}.

\subsection{Partitioning and mesh}

Meshing a plate with two circular patches bonded to it is the first practical obstacle. A free mesh degrades
badly where it has to transition around a circular footprint, and the elements it produces there are distorted
enough to reflect the wave on their own. The plate was therefore partitioned before meshing, both around the
transducer footprints and around the crack, so the mesh could be built from mappable regions
(Figures~\ref{fig:p2geom} and~\ref{fig:p2mesh}). Edge seeding was then used to refine the mesh locally near
the discs and near the crack, where the wavelength of the scattered field is shortest.

\begin{figure}[H]
\centering
\begin{minipage}[t]{0.48\textwidth}
  \centering
  \includegraphics[width=\textwidth]{p2_plate_healthy.png}
\end{minipage}\hfill
\begin{minipage}[t]{0.48\textwidth}
  \centering
  \includegraphics[width=\textwidth]{p2_plate_cracked.png}
\end{minipage}
\caption{Plate partitioning before meshing. Left: healthy plate, with the two circular transducer footprints
partitioned into mappable faces. Right: the same plate with the 25\,mm crack partition added at the centre.}
\label{fig:p2geom}
\end{figure}

\begin{figure}[H]
\centering
\begin{minipage}[t]{0.48\textwidth}
  \centering
  \includegraphics[width=\textwidth]{p2_pzt_pair.png}\\[2pt]
  {\small (a) The two PZT discs, partitioned}
\end{minipage}\hfill
\begin{minipage}[t]{0.48\textwidth}
  \centering
  \includegraphics[width=\textwidth]{p2_mesh_disc.png}\\[2pt]
  {\small (b) Meshed disc}
\end{minipage}

\vspace{8pt}

\begin{minipage}[t]{0.48\textwidth}
  \centering
  \includegraphics[width=\textwidth]{p2_mesh_pzt_seat.png}\\[2pt]
  {\small (c) Plate mesh under a transducer}
\end{minipage}\hfill
\begin{minipage}[t]{0.48\textwidth}
  \centering
  \includegraphics[width=\textwidth]{p2_mesh_crack.png}\\[2pt]
  {\small (d) Plate mesh around the crack}
\end{minipage}
\caption{Discretisation of the two-transducer model. The mesh is refined locally where the geometry changes
and where the scattered field is expected to have the shortest wavelength.}
\label{fig:p2mesh}
\end{figure}

\subsection{Loads, boundary conditions and solution strategy}

No mechanical load is applied to the plate or to the discs. The excitation is entirely electrical. The bonded
face of each disc, where it meets the plate, is held at zero volts, and the free face of the actuator is driven
with a $200$\,kHz sinusoidal potential of $10$\,V amplitude. The free face of the sensor is left free to develop
a potential, and that potential is the output of the analysis.

The solver strategy is dictated by an Abaqus restriction. The fully coupled piezoelectric continuum element
is available in Abaqus/Standard but not in Abaqus/Explicit, whereas explicit time integration is what makes
a wave-propagation analysis affordable. The two requirements were reconciled by splitting the problem: the
aluminium plate is solved with Dynamic, Explicit, the PZT discs with Dynamic, Implicit, and the two analyses
are run simultaneously and exchange data at their common interface through an Abaqus co-execution
(Figure~\ref{fig:p2coexec}). Getting the interface right, so that the partitioned faces on the underside of each
disc and on the surface of the plate coincide, is the part of the setup that most often prevents the analysis
from running.

\begin{figure}[H]
\centering
\includegraphics[width=0.62\textwidth]{p2_coexecution.png}
\caption{Co-execution manager. The healthy and cracked cases are each solved as a coupled
Abaqus/Standard--Abaqus/Explicit pair, both completing successfully.}
\label{fig:p2coexec}
\end{figure}

Two analyses were run, identical except for the crack: a healthy plate and a plate in which the crack was
introduced as a cut extruded through the thickness. Both were run to $80$\,\textmu s.

\subsection{Damage indices considered}

Once a healthy signal $S_{r}(i)$ and a current signal $S_{d}(i)$ are available, a scalar damage index has to be
extracted from them. Several were reviewed before two were selected.

\paragraph{Signal difference coefficient (SDC).}
Defined as one minus the normalised cross-correlation of the two signals,

\begin{equation}
  \mathrm{SDC} = 1 - \frac{\sum_{i=1}^{N}\tilde{S}_{d}(i)\,\tilde{S}_{r}(i)}
  {\sqrt{\sum_{i=1}^{N}\tilde{S}_{d}^{2}(i)\ \sum_{i=1}^{N}\tilde{S}_{r}^{2}(i)}}
  \label{eq:sdc}
\end{equation}

where $\tilde{S}(i) = S(i)-\bar{S}$ is the mean-removed signal and the sums run over the $N$ samples of the
record. Identical signals give $\mathrm{SDC}=0$ and
completely dissimilar signals give $\mathrm{SDC}=1$. Because it uses the full waveform, SDC responds to
changes in shape and phase as well as amplitude.

\paragraph{Normalised root-mean-square deviation.}

\begin{equation}
  \mathrm{NRMSD} = \sqrt{\frac{\sum_{i=1}^{N}\left[S_{d}(i)-S_{r}(i)\right]^{2}}{\sum_{i=1}^{N}S_{r}(i)^{2}}}
\end{equation}

which measures the energy of the difference between the signals, normalised by the energy of the baseline.
It has a clear physical reading as a fraction of signal energy changed, and thresholds around $0.2$ are
commonly quoted as the level above which a change exceeds normal environmental variation.

\paragraph{Peak amplitude ratio.}
The simplest index of all,

\begin{equation}
  A_{\text{ratio}} = \frac{A_{\text{damaged}}}{A_{\text{baseline}}}
\end{equation}

comparing the peak amplitude of the first arriving packet. Scattering at the crack removes energy from the
direct path, so a ratio below one indicates attenuation.

\paragraph{Other features.}
Time-of-flight shift, residual signal energy, and time-frequency features from the continuous wavelet
transform, the short-time Fourier transform or the Hilbert envelope were also reviewed. The wavelet-energy
feature is the one that becomes central in Stage 2, and it is defined in Section~\ref{sec:di3}.

\subsection{Results}
\label{sec:stage1results}

Figure~\ref{fig:p2wavefield} shows the wave leaving the actuator and spreading across the plate, confirming
that the electrical excitation does generate a propagating elastic field. Figure~\ref{fig:p2signals} shows the
electric potential recorded at the sensor electrode in each of the two states.

\begin{figure}[H]
\centering
\includegraphics[width=0.78\textwidth]{p2_wavefield.png}
\caption{Von Mises stress at $16$\,\textmu s, showing the wave radiating from the actuator across the healthy
plate.}
\label{fig:p2wavefield}
\end{figure}

\begin{figure}[H]
\centering
\includegraphics[width=0.82\textwidth]{p2_signal_healthy.png}\\[2pt]
{\small (a) Healthy plate}\\[8pt]
\includegraphics[width=0.82\textwidth]{p2_signal_cracked.png}\\[2pt]
{\small (b) Cracked plate}
\caption{Electric potential at the sensor electrode over $80$\,\textmu s. Both traces are flat until about
$27$\,\textmu s, which agrees with the $26.2$\,\textmu s predicted from the $S_{0}$ group velocity.}
\label{fig:p2signals}
\end{figure}

Both signals stay at zero until roughly $27$\,\textmu s and then begin to oscillate. That arrival time is the
first useful check on the model: the $S_{0}$ group velocity from Table~\ref{tab:dispersion} predicts
$26.2$\,\textmu s over the $139.7$\,mm path, so the wave in the model is travelling at the speed the
analytical dispersion relation says it should.

Overlaying the two traces (Figure~\ref{fig:p2overlay}) shows the effect of the crack. The cracked-state
signal is lower in amplitude throughout, and the two waveforms drift progressively out of phase as time
increases.

\begin{figure}[H]
\centering
\includegraphics[width=0.85\textwidth]{p2_signal_overlay.png}
\caption{Sensor signal from the healthy plate and from the cracked plate, overlaid. The crack reduces the
amplitude and shifts the phase of the later part of the record.}
\label{fig:p2overlay}
\end{figure}

The two selected indices give the values in Table~\ref{tab:res2}.

\begin{table}[H]
\centering
\caption{Damage indices for the two-transducer model.}
\begin{tabular}{l c}
\toprule
Quantity & Value \\
\midrule
Signal difference coefficient & $0.989$ \\
Peak amplitude, healthy       & $0.195$ \\
Peak amplitude, cracked       & $0.157$ \\
Amplitude drop                & $19.5\%$ \\
\bottomrule
\end{tabular}
\label{tab:res2}
\end{table}

The amplitude drop of $19.5\%$ is unambiguous. A change of 5--10\% is normally taken as suggestive of
damage and anything above 10\% as a strong indication, so a drop of nearly 20\% is a clear detection. The
SDC of $0.989$ is also a clear detection, but its magnitude needs care in interpretation. Equation
\eqref{eq:sdc} is evaluated at zero lag over the whole $80$\,\textmu s record, and because the two waveforms
drift out of phase towards the end of that window their zero-lag correlation collapses towards zero, which
sends SDC towards one. The index has effectively saturated. It confirms that the signals differ, but at that
value it can no longer rank one difference against another, which matters as soon as several sensing paths
have to be compared.

\subsection{From detection to localisation}

The two-transducer model was built to answer one question, and by design it carries no information for a
second one. Three specific requirements separate what it does from what localisation needs, and Stage 2 is
built around meeting all three.

The first is geometric. A single actuator-sensor pair yields one number, and one number cannot locate
anything. The crack in this model happens to sit on the direct path, so it produces a large index; a crack of
the same size sitting 30\,mm to the side would produce a smaller index, and the two situations are
indistinguishable from a single measurement. A single path can therefore only detect, and reaching Level~2
of Table~\ref{tab:levels} requires several paths crossing the monitored area at different angles.

The second is in the excitation. Driving the actuator with a continuous 200\,kHz sinusoid, as was done here,
means the plate is being fed energy for the entire analysis. The received signal grows steadily throughout
the record rather than arriving as a discrete packet, and by $80$\,\textmu s it also contains reflections from
the plate edges superposed on the direct arrival. That is workable when the only question is whether the
signal changed, and it is unusable when the arrival time of a specific wave packet has to be measured. A
windowed tone burst is required instead, and that is what Stage 2 uses.

The third is a modelling issue. The crack here was created by cutting material out of the plate, which
means the healthy and cracked models have different meshes. Part of the difference between the two signals
is therefore attributable to the change in discretisation rather than to the crack. For a detection threshold
this is tolerable, and for a quantitative per-path index that has to be compared across six paths it is not.

%-----------------------------------------------------
\section{Stage 2: crack localisation with a four-transducer network}
\label{sec:stage2}
%-----------------------------------------------------

The Stage 2 model extends the same finite-element approach to meet all three requirements above.

\subsection{Layout and geometry}

Four PZT discs are placed at the corners of a square on a $400\times400\times1.6$\,mm aluminium plate. Four
transducers give $\binom{4}{2}=6$ unique sensing paths, four along the sides and two along the diagonals.
The crack is placed off-centre, on the P1--P3 diagonal, so that one path passes directly through it while the
remaining five pass at increasing perpendicular offsets. Table~\ref{tab:geom3} gives the layout and
Figure~\ref{fig:p3layout} shows it.

\begin{table}[H]
\centering
\caption{Geometry of the four-transducer model.}
\fitwidth{%
\begin{tabular}{l l}
\toprule
Plate & $400 \times 400 \times 1.6$\,mm, aluminium \\
PZT discs & 4, diameter $12.7$\,mm, thickness $0.254$\,mm \\
PZT positions & P1$(-75,-75)$, P2$(75,-75)$, P3$(75,75)$, P4$(-75,75)$\,mm \\
Crack & $25$\,mm long, through-thickness, centred at $(-35,-35)$\,mm \\
Path lengths & 4 sides at $150$\,mm, 2 diagonals at $212.1$\,mm \\
\bottomrule
\end{tabular}}
\label{tab:geom3}
\end{table}

\begin{figure}[H]
\centering
\includegraphics[width=0.86\textwidth]{p3_plate_layout.png}
\caption{Plate part seen from above. Each of the four transducer footprints carries the butterfly partition
inside its $30\times30$\,mm structured box, and the small partition left of centre bounds the crack.}
\label{fig:p3layout}
\end{figure}

The plate size follows from a timing constraint. This is the smallest plate that keeps the first edge reflection
outside the $S_{0}$ arrival window on all six paths: reflections arrive at $54.7$\,\textmu s, while the latest
$S_{0}$ window closes at $48.1$\,\textmu s. That removes the edge-reflection contamination which limited the
Stage 1 record.

\subsection{Crack representation}

Instead of cutting material away, the crack is modelled as a zero-thickness seam. A seam duplicates the nodes
on the crack face so the two faces can separate, while leaving the rest of the mesh untouched. The healthy and
cracked models therefore have identical meshes, element for element: both contain $1{,}040{,}892$ elements,
and the seam adds exactly 150 duplicated nodes on the crack face. The only difference between the two
analyses is the crack itself, which is what a per-path damage index requires if it is to mean anything. A notch
variant with a 1\,mm wide cut was also built in the same CAE model, giving a second crack representation to
compare the seam against; its co-simulations were not run for this report.

\subsection{Mesh design from the dispersion relations}

The element size comes straight out of Table~\ref{tab:dispersion}. The shorter of the two wavelengths at
200\,kHz belongs to $A_{0}$ at $7.90$\,mm, and the usual requirement is ten elements per wavelength, giving
an in-plane element size of $0.8$\,mm, with four elements through the thickness.

Meshing the plate with four circular footprints raised the same difficulty as in Stage 1, in a more severe
form. A free mesh produced 224 elements with zero or negative volume and 952 elements with an aspect ratio
above 100, and it also introduced a spurious reflection at the mesh transition that was visible in the signals.
The fix was to partition each footprint into a butterfly pattern of five mappable faces, an inner square plus
four radial segments reaching the corners of a $30\times30$\,mm structured box, and to apply the same
partition to the discs themselves (Figure~\ref{fig:p3disc}). The plate away from the boxes was then meshed
with a structured technique and only the boxes with sweep or advancing-front. The result is a fully
structured hexahedral mesh of C3D8E elements with zero failed and zero warning elements, and the spurious
reflection disappeared with it.

\begin{figure}[H]
\centering
\begin{minipage}[t]{0.48\textwidth}
  \centering
  \includegraphics[width=\textwidth]{p3_disc_geometry.png}\\[2pt]
  {\small (a) Meshed PZT disc}
\end{minipage}\hfill
\begin{minipage}[t]{0.48\textwidth}
  \centering
  \includegraphics[width=\textwidth]{p3_butterfly.png}\\[2pt]
  {\small (b) Plate mesh beneath transducer P1}
\end{minipage}
\caption{Butterfly partition into an inner square and four radial segments, applied to the disc itself (a) and
to the plate under it (b). The inner square meshes as a regular grid, and the four segments carry the
transition out to the edge of the $30\times30$\,mm structured box. That box edge is the visible boundary in
(b) between the swept region and the uniform 0.8\,mm grid outside it.}
\label{fig:p3disc}
\end{figure}

A seam has zero thickness, so the mesh closes over it and the crack leaves nothing to see from above.
Figure~\ref{fig:p3crack}(a) therefore marks the seam on the surrounding mesh, while
Figure~\ref{fig:p3crack}(b) shows the notch variant, where the 1\,mm cut is a real hole in the mesh.

\begin{figure}[H]
\centering
\begin{minipage}[t]{0.48\textwidth}
  \centering
  \includegraphics[width=\textwidth]{p3_crack_mesh.png}\\[2pt]
  {\small (a) Seam model, crack face marked}
\end{minipage}\hfill
\begin{minipage}[t]{0.48\textwidth}
  \centering
  \includegraphics[width=\textwidth]{p3_notch_mesh.png}\\[2pt]
  {\small (b) Notch variant, 1\,mm cut}
\end{minipage}
\caption{The crack region. In (a) the uniform 0.8\,mm mesh runs through the crack location undisturbed, with
the seam drawn in red along the line where its faces are duplicated. In (b) the notch variant removes
material, so the mesh wraps around the cut and the element count rises to $1{,}231{,}288$.}
\label{fig:p3crack}
\end{figure}

\subsection{Excitation}

The continuous sinusoid of Stage 1 is replaced by a four-cycle Hann-windowed tone burst,

\begin{equation}
  V(t) = \frac{V_{0}}{2}\left[1-\cos\!\left(\frac{2\pi f_{0} t}{N_{\text{cyc}}}\right)\right]
         \sin\!\left(2\pi f_{0}t\right),
  \qquad 0 \le t \le \frac{N_{\text{cyc}}}{f_{0}}
\end{equation}

with $V_{0}=10$\,V, $f_{0}=200$\,kHz and $N_{\text{cyc}}=4$. The window narrows the bandwidth, which
limits dispersion of the packet, and it produces a wave packet with a definite start and end, which is what
makes an arrival time measurable.

\subsection{Electrode modelling and co-simulation}

The bottom electrode of every disc is grounded. The actuator's top electrode is driven with the tone burst.
For the sensors, the top electrode is tied to a single equipotential node through \texttt{*EQUATION}
constraints on the electric-potential degree of freedom, which reproduces the behaviour of a real continuous
metal electrode instead of treating every finite-element node on the surface as an independent voltage
output (Figure~\ref{fig:p3bc}). The electric potential at each electrode's master node is what is recorded.

\begin{figure}[H]
\centering
\includegraphics[width=0.55\textwidth]{p3_electrical_bc.png}
\caption{Electrical boundary conditions on one transducer: grounded bottom electrode, driven or sensing top
electrode.}
\label{fig:p3bc}
\end{figure}

The plate is again solved in Abaqus/Explicit and the discs in Abaqus/Standard, coupled through a
surface-based co-simulation with lockstep incrementation. One co-simulation is required per actuator per
plate state, so eight runs were performed in total, covering actuators P1 to P4 in both healthy and cracked
states. Three actuators are enough to cover all six unique undirected paths; the fourth actuator repeats those
paths in the reverse direction and was used as an independent reciprocity check.

\subsection{Verification}

The following checks were performed before any result was interpreted.

\begin{itemize}[leftmargin=*,itemsep=2pt]
  \item Healthy and cracked meshes are identical element for element, $1{,}040{,}892$ elements each, with the
  seam contributing exactly 150 duplicated nodes.
  \item Mesh quality verification with analysis checks reports zero failed and zero warning elements on all
  plate variants.
  \item A data check on the plate model and on each disc model completes with zero errors. This caught two
  real modelling errors during development, a duplicate equation-constraint conflict created by the butterfly
  partition and a missing local material orientation on the anisotropic piezoelectric material, both fixed
  before the production runs.
  \item All eight co-simulations reached completed status.
\end{itemize}

\subsection{Damage index and probability-based diagnostic imaging}
\label{sec:di3}

Following Nasiri et al.~\cite{Nasiri2022}, the damage-sensitive feature is the wavelet energy of the signal
inside a time window,

\begin{equation}
  F = \sum_{t}\sum_{f}\left|W(t,f)\right|^{2}
\end{equation}

where $W(t,f)$ is the continuous wavelet transform of the signal using a complex Morlet wavelet. The
damage index for a path is the relative change in that energy,

\begin{equation}
  \mathrm{DI} = \frac{\left|F_{\text{damaged}} - F_{\text{healthy}}\right|}{F_{\text{healthy}}}.
\end{equation}

A wavelet-energy feature was preferred over the SDC used in Stage 1 for the reason given in
Section~\ref{sec:stage1results}: SDC saturates, and a per-path index that has to be compared across six
paths must stay graded.

Turning six scalar indices into a two-dimensional map is the job of the PDI algorithm. Each path contributes
an elliptical weight function centred on the straight line joining its two transducers,

\begin{align}
  r(x,y) &= \frac{d_{i}+d_{j}}{d_{ij}} - 1 \\[2pt]
  \gamma &= \frac{V_{S_{0}}\,T_{S_{0}}}{d_{ij}} \\[2pt]
  \beta  &= \frac{\sqrt{2}}{n}\,\gamma, \qquad n = 4 \\[2pt]
  W(x,y) &= \exp\!\left[-\left(\frac{r}{\beta}\right)^{2}\right] \\[2pt]
  P(x,y) &= \sum_{\text{paths}} \mathrm{DI}_{ij}\,W_{ij}(x,y)
\end{align}

where $d_{i}$ and $d_{j}$ are the distances from the point $(x,y)$ to the two transducers, $d_{ij}$ is the
direct path length, and $T_{S_{0}} = N_{\text{cyc}}/f_{0}$ is the duration of the $S_{0}$ packet. The
quantity $r$ is zero on the direct path and grows away from it, so each path produces a band of high weight
that is widest at the midpoint and narrows towards the transducers. Summing the weighted indices over all
paths concentrates probability where the bands from the most strongly affected paths overlap. The
coefficient $\beta$ controls how wide those bands are. Tying $\beta$ to the physical length of the wave packet,
which is what \cite{Nasiri2022} contributes, replaces the trial-and-error value used in earlier formulations.

\subsection{Results}

\subsubsection{Sensor signals}

Figure~\ref{fig:p3signals} compares the received signal on all six paths, healthy against cracked, ordered by
how close each path passes to the crack. The vertical marker on each subplot is the $S_{0}$ arrival time
predicted from the group velocity of Table~\ref{tab:dispersion}, and it lands on the observed onset in every
case.

\begin{figure}[H]
\centering
\includegraphics[height=0.78\textheight]{p3_signal_comparison.png}
\caption{$S_{0}$ wave packet on all six sensing paths, healthy plate against cracked plate. The predicted
$S_{0}$ arrival time is marked on each subplot.}
\label{fig:p3signals}
\end{figure}

The pattern across the six paths follows the crack geometry (Figure~\ref{fig:p3drop} and
Table~\ref{tab:di3}). The P1--P3 diagonal, which passes straight through the crack at zero offset, loses
$23.4\%$ of its peak amplitude and changes shape after the first arrival. The three paths offset by
40--50\,mm change by a few percent, in either direction, which is consistent with the scattered field
interfering constructively or destructively at the sensor depending on the path geometry. The two paths
110\,mm from the crack show no change at all, to four decimal places in the damage index. That last
observation also confirms the healthy and cracked models are identical apart from the crack, since meshes
that differed would have registered something on those distant paths.

\begin{figure}[H]
\centering
\includegraphics[width=0.85\textwidth]{p3_amplitude_drop.png}
\caption{Peak amplitude drop from healthy to cracked state for each sensing path, annotated with the path's
perpendicular offset from the crack.}
\label{fig:p3drop}
\end{figure}

\begin{table}[H]
\centering
\caption{Per-path damage indices and amplitude drops. A negative drop means the amplitude increased.}
\fitwidth{%
\begin{tabular}{l c c c c}
\toprule
Path & Length (mm) & Offset from crack (mm) & Damage index & Amplitude drop (\%) \\
\midrule
P1--P3 & 212.1 & 0     & 0.4065 & 23.4 \\
P1--P2 & 150.0 & 40    & 0.1400 & $-7.5$ \\
P1--P4 & 150.0 & 40    & 0.0696 & 3.5 \\
P2--P4 & 212.1 & 49.5  & 0.0996 & $-4.8$ \\
P2--P3 & 150.0 & 110   & 0.0001 & 0.0 \\
P3--P4 & 150.0 & 110   & 0.0002 & 0.0 \\
\bottomrule
\end{tabular}}
\label{tab:di3}
\end{table}

\subsubsection{Localisation by probability-based diagnostic imaging}

Feeding the six damage indices into the PDI algorithm produces the map in Figure~\ref{fig:p3pdi}. The peak of
the map is at $(-35,-45)$\,mm against a true crack location of $(-35,-35)$\,mm, an error of $10.0$\,mm. On a
plate spanning 400\,mm that is an error of about $2.5\%$, obtained from four transducers and six paths.

\begin{figure}[H]
\centering
\includegraphics[width=0.72\textwidth]{p3_pdi_map.png}
\caption{Damage-probability map from the PDI algorithm with Gaussian elliptical weighting. Green cross:
true crack location. Cyan cross: the peak of the map, which is the algorithm's estimate.}
\label{fig:p3pdi}
\end{figure}

\subsubsection{Localisation by time-of-flight delay-and-sum}

A second and independent estimator was implemented on the same extracted signals, with no further
simulation. Subtracting the healthy signal from the cracked signal on each path,

\begin{equation}
  v_{\text{residual}}(t) = v_{\text{cracked}}(t) - v_{\text{healthy}}(t),
\end{equation}

removes the direct arrival, which is common to both states, and leaves the field scattered by the crack. For
every candidate point on the plate the expected time of flight from actuator to that point to sensor is
computed, the residual envelope is sampled at that delay, and the contributions from all six paths are
summed. This is a migration image, the same principle as the total focusing method used in phased-array
ultrasonic inspection.

\begin{figure}[H]
\centering
\includegraphics[width=0.72\textwidth]{p3_tof_map.png}
\caption{Delay-and-sum image formed from the time of flight of the crack-scattered wave on all six paths.}
\label{fig:p3tof}
\end{figure}

The estimate is $(-37,-110)$\,mm, an error of $75.0$\,mm, which is much worse than the PDI result. Two
effects account for this. The four-cycle tone burst lasts about $20$\,\textmu s, which at the $S_{0}$ group
velocity corresponds to a wave packet roughly 100\,mm long in space, so the envelope peak used to measure
the arrival is blurred over tens of millimetres whatever method reads it. On top of that, four transducers at
the corners give six paths with limited angular diversity, and the image cannot distinguish the true crossing
point from other places where several path ellipses happen to intersect. The bright ring away from the crack
in Figure~\ref{fig:p3tof} is that ghosting, a known behaviour of sparse-array delay-and-sum imaging. The
full-field check below establishes that the underlying model and signals are sound.

\subsubsection{Full-field validation}

The last check verifies the model itself. It is not deployable, since it requires the entire simulated
displacement field, which no realistic sensor network can measure. The whole-plate displacement
field was extracted from the explicit output database at $t=40$\,\textmu s for both the healthy and cracked
models under actuator P1, and the two fields were subtracted node by node.

\begin{figure}[H]
\centering
\includegraphics[width=0.72\textwidth]{p3_wavefield_diff.png}
\caption{$\left|U_{\text{crack}} - U_{\text{healthy}}\right|$ over the whole plate at $t=40$\,\textmu s. The
circular wavefront is centred on the true crack location.}
\label{fig:p3field}
\end{figure}

The difference field is a clean circular wavefront expanding from a point that coincides with the true crack
location to within the 2\,mm resolution of the grid used. The crack behaves as a secondary point source
exactly where it was placed. What the point-sensor methods infer indirectly from six voltage traces is
therefore visible directly here, and it puts an upper bound on how much of the $10$\,mm PDI error the model
itself could account for: the model places the crack correctly, so the remaining error is set by what four
sensors and six paths can resolve.

\subsection{Comparison of the three estimators}

\begin{table}[H]
\centering
\caption{Localisation error of the three methods applied to the same simulation.}
\fitwidth{%
\begin{tabular}{l l c}
\toprule
Method & Input used & Localisation error \\
\midrule
Classical PDI (wavelet energy) & 6 paths + healthy baseline & $10.0$\,mm \\
Time-of-flight delay-and-sum & 6 paths + healthy baseline & $75.0$\,mm \\
Full-field wavefield difference & full displacement field & $\sim 2$\,mm (validation only) \\
\bottomrule
\end{tabular}}
\label{tab:methods3}
\end{table}

Three points follow from Table~\ref{tab:methods3}.

The energy change in a guided-wave signal correlates strongly with how close the direct path passes to the
crack. The path that intersects the crack changes by far the most, and paths 110\,mm away do not change at
all. Elliptical weighting turns that correlation into a usable image, and it is the best-performing
sensor-based method tested here.

Changing the feature from amplitude to time of flight did not improve accuracy with this layout. With four
corner transducers the aperture is too sparse to separate the true crossing point from ghost intersections,
and the finite duration of the tone burst blurs any envelope-based time measurement regardless of which
feature is used. The elliptical Gaussian weight in the PDI algorithm already absorbs that blur, which is why
it tolerates the sparse layout better.

The full-field check settles where the residual error comes from. The crack sits where it was placed in the
model, so the $10$\,mm is what this sensor layout can resolve.

\subsection{Comparison against published PDI results}
\label{sec:litcompare}

The PDI algorithm used here is the one developed by Nasiri et al.~\cite{Nasiri2022} and applied to sensor-layout
design by Keshavarz Motamed et al.~\cite{Motamed2019}, and both report actual localisation errors that this
result can be checked against.

Nasiri et al.\ instrument a $500\times500\times1$\,mm stiffener-reinforced aluminium plate with nine PZTs (56
sensing paths) and report errors of $2.2$--$41.9$\,mm across three single-crack scenarios with their best weight
function ($\beta$-GWD), averaging about $22$\,mm; their fatigue cracks are $4$\,mm long. Keshavarz Motamed
et al.\ use the same $400\times400\times1$\,mm aluminium plate geometry as this report, with a $200\times200$\,mm
central monitoring area and a $20$\,mm crack, and report two relevant results: with 8 sensors, finite-element
verification of three single-crack cases gives errors of $7.8$--$19.5$\,mm; and, across a statistical sweep of 121
crack positions, a circular array needs 14--20 sensors to keep 80\% of predictions within $10$\,mm.

\begin{table}[H]
\centering
\caption{Published PDI localisation errors compared against the result of Section~\ref{sec:stage2}.}
\fitwidth{%
\begin{tabular}{l c c c c}
\toprule
Source & Plate & Crack & Sensors & Localisation error \\
\midrule
Nasiri et al.~\cite{Nasiri2022} & $500\times500$\,mm, stiffened & $4$\,mm & 9 (56 paths) & $2.2$--$41.9$\,mm (3 cases) \\
Motamed et al.~\cite{Motamed2019} & $400\times400$\,mm, plain & $20$\,mm & 8 & $7.8$--$19.5$\,mm (3 cases) \\
Motamed et al.~\cite{Motamed2019} & $400\times400$\,mm, plain & $20$\,mm & 14--20 & $10$\,mm at 80th pctl. \\
\midrule
This report (Section~\ref{sec:stage2}) & $400\times400$\,mm, plain & $25$\,mm & 4 (6 paths) & $10.0$\,mm (1 case) \\
\bottomrule
\end{tabular}}
\label{tab:litcompare}
\end{table}

The comparison has to be read with its caveats attached, and the crack itself is the largest one: at $25$\,mm,
the crack in this report is roughly six times longer than the $4$\,mm fatigue cracks in Nasiri et al.\ and about a
quarter longer than the $20$\,mm crack in Keshavarz Motamed et al. A longer crack scatters more of the wave
and is intrinsically easier to detect and locate, so part of the gap in sensor count is the comparison rewarding
a larger flaw rather than a better method. The crack position adds a second advantage: it sits exactly on the
P1--P3 path (Table~\ref{tab:di3}), which is the single most favourable geometry for any of these arrays, whereas
the published numbers are statistics over many crack positions and orientations, most of which are not that
well aligned. Nasiri et al.'s plate also carries a stiffener and two rows of rivet holes, which give the wave field
more to scatter from than the plain plate used here. Taken together, this report's $10.0$\,mm is one data point
from a favourable case, not a distribution, and it should be read as evidence that the PDI method can reach
this accuracy with four sensors under good conditions, not as evidence that four sensors are sufficient for
this plate in general. Within that reading, the result is still a genuine one: a bachelor's-level model with a
sparse, unoptimised four-corner layout reaches the same accuracy that the cited graduate-level studies, from
the same research group, needed considerably denser arrays to reach even on their easiest cases.

%-----------------------------------------------------
\section{Discussion}
%-----------------------------------------------------

Read together, the three phases form a single argument: choose the method, demonstrate that it detects,
then extend it until it localises.

The method review fixed the direction of the work. Once the requirement was stated as high sensitivity to a
small crack, a large monitored area per sensor, and permanent installation, guided-wave monitoring with
bonded piezoelectric transducers was the only family that satisfied all three. The review supplies the
argument for that choice.

The two-transducer model established that the physics could be captured numerically. The arrival time
matched the analytical $S_{0}$ group velocity, the crack produced a $19.5\%$ amplitude drop, and the
implicit-explicit co-execution proved to be a workable route around the absence of coupled piezoelectric
elements in Abaqus/Explicit. It also made clear three specific requirements a localisation-capable model
would have to satisfy: multiple paths, since a single path cannot localise; a windowed excitation, since a
continuous tone prevents any time-of-flight measurement and admits edge reflections; and a crack
representation that leaves the mesh unchanged, since a crack cut out of the mesh contaminates the
comparison with a discretisation difference.

The four-transducer model meets all three. Six paths replaced one, a windowed tone burst replaced the
continuous sinusoid, and a seam replaced the cut so that the healthy and cracked meshes stayed identical.
That last change buys a piece of verification the Stage 1 model could not offer: with the two meshes identical
element for element, the damage index of $0.0001$ reported by the two distant paths measures the numerical
noise floor of the whole procedure. Against that floor, the $0.4065$ index on the path through the crack is
unambiguous.

The gap between the PDI result and the delay-and-sum result is what the comparison teaches. Both estimators
consume the same six signals, and one is seven times more accurate than the other, so the difference lies in
how much each method demands of the data. Delay-and-sum needs a resolvable arrival time from every path,
which a 100\,mm-long wave packet on a sparse array cannot supply. PDI needs only a graded scalar per path,
which the same signals do supply. Array density, then, is what limits accuracy here, and switching damage
feature will not move it much.

%-----------------------------------------------------
\section{Conclusions}
%-----------------------------------------------------

\begin{enumerate}[leftmargin=*,itemsep=3pt]
  \item A comparison of SHM method families against the requirements of a thin metallic panel selects
  guided-wave monitoring with bonded piezoelectric transducers, which is the only approach combining
  sensitivity to small cracks with a large monitored area per sensor and permanent installation.

  \item A coupled finite-element model of an aluminium plate with two bonded PZT discs, solved as an
  implicit-explicit co-execution in Abaqus, reproduces $S_{0}$ Lamb-wave propagation at the velocity
  predicted by the Rayleigh--Lamb equations. A 25\,mm crack on the direct path produces a $19.5\%$ drop in
  peak amplitude and a signal difference coefficient of $0.989$, which is a clear Level~1 detection.

  \item The same modelling approach extended to four transducers gives six independent sensing paths.
  Damage indices computed from the continuous wavelet energy of each path follow the expected physical
  ordering, from $0.4065$ on the path through the crack down to $0.0001$ on the paths 110\,mm away.

  \item Combining those indices through a probability-based diagnostic imaging algorithm with elliptical
  Gaussian weighting locates the crack to within $10$\,mm of its true position on a 400\,mm plate, which is
  a Level~2 diagnosis from four sensors. That error sits inside the $2.2$--$41.9$\,mm range reported for the
  same class of algorithm with denser arrays in two published studies (Section~\ref{sec:litcompare}), one of
  which uses the identical $400\times400$\,mm plate and needs 14--20 sensors for a comparable error at its
  80th percentile, against the four used here on a single, favourably positioned case with a longer crack.

  \item A time-of-flight delay-and-sum image built from the same signals gives a $75$\,mm error, caused by
  the sparse aperture of a four-corner array together with the spatial length of the tone burst.

  \item A full-field wavefield-difference check shows the crack acting as a secondary point source at its
  true location to within the 2\,mm grid resolution. The residual localisation error is therefore set by what
  the sensor layout can resolve.
\end{enumerate}

%-----------------------------------------------------
\section{Future work}
%-----------------------------------------------------

\paragraph{Denser and better-distributed arrays.}
Table~\ref{tab:methods3} points to array sparsity as the factor limiting accuracy. Six, eight or nine
transducers distributed around the plate, instead of four at the corners, would increase both the number of
independent paths and the angular diversity. The delay-and-sum estimator, whose ghosting is a direct
consequence of too few paths, would gain the most.

\paragraph{Feature fusion.}
Amplitude-based and time-of-flight features fail in different ways, the first through interference of the
scattered field and the second through sparse-aperture ghosting. Combining them into a single fused image
would prevent either failure mode from dominating.

\paragraph{Baseline-free operation.}
Every index used here needs a recorded healthy signal, and in service that baseline drifts with temperature
and loading. Instantaneous-baseline methods, which compare geometrically identical paths within the same
measurement instead of comparing against a stored record, remove that dependence
\cite{Alem2016,Alem2017,Alem2021}, and they are the natural next step for making this approach usable
outside a simulation.

\paragraph{Optimal sensor placement.}
The layout can be treated as a design variable and optimised against localisation accuracy
\cite{Motamed2019,Hassani2023Opt}, which turns the trial layout used here into a designed one.

\paragraph{Machine learning on the signal database.}
The finite-element model can generate labelled signals for many crack positions, lengths and orientations at a
cost that physical experiments cannot match. That database is what a learned classifier or regressor needs,
and it would let damage type and severity, Levels 3 and 5 of Table~\ref{tab:levels}, be estimated alongside
position \cite{Jia2023,YangZ2023}.
